% =========================================================
% Computational Linear Algebra — Determinants and the Cross Product
% (Strang §5.2 / Johnston §1.A / Ashraf §1.4)
% =========================================================

\section{Determinants and the Cross Product}

\begin{tcolorbox}[colback=gray!6, colframe=gray!40, arc=2pt,
  left=6pt, right=6pt, top=4pt, bottom=4pt,
  title={\small\textbf{Determinants and Cross Product --- Quick Reference}},
  fonttitle=\small\bfseries]
\begin{tabular}{@{}p{0.30\textwidth}p{0.63\textwidth}@{}}
\textbf{$2\times 2$ det}  & $\det\begin{pmatrix} a & b \\ c & d \end{pmatrix} = ad - bc$. \\[6pt]
\textbf{$3\times 3$ det}  & Cofactor expansion along any row or column. \\[4pt]
\textbf{Cross product}    & $\mathbf{v} \times \mathbf{w}$: vector orthogonal to both $\mathbf{v}$ and $\mathbf{w}$ in $\mathbb{R}^3$, computed via $3\times 3$ determinant mnemonic. \\[4pt]
\textbf{Magnitude}        & $\|\mathbf{v} \times \mathbf{w}\| = \|\mathbf{v}\|\,\|\mathbf{w}\|\sin\theta$; equals area of parallelogram. \\[4pt]
\textbf{Direction}        & Right-hand rule; $\mathbf{v} \times \mathbf{w} = -(\mathbf{w} \times \mathbf{v})$. \\
\end{tabular}
\end{tcolorbox}

% ---------------------------------------------------------
\subsection{The $3 \times 3$ Determinant}

\begin{tcolorbox}[colback=propbox, colframe=propborder, arc=2pt,
  left=6pt, right=6pt, top=4pt, bottom=4pt,
  title={\small\textbf{Definition ($3 \times 3$ Determinant via Cofactor Expansion)}},
  fonttitle=\small\bfseries]
For $A \in \mathbb{R}^{3 \times 3}$, expanding along the first row:
\[
  \det(A) = a_{11}(a_{22}a_{33} - a_{23}a_{32})
           - a_{12}(a_{21}a_{33} - a_{23}a_{31})
           + a_{13}(a_{21}a_{32} - a_{22}a_{31}).
\]
The signs alternate: $+, -, +$ along the first row.
\end{tcolorbox}

% ---------------------------------------------------------
\subsection{The Cross Product}

\begin{tcolorbox}[colback=propbox, colframe=propborder, arc=2pt,
  left=6pt, right=6pt, top=4pt, bottom=4pt,
  title={\small\textbf{Definition (Cross Product)}},
  fonttitle=\small\bfseries]
For $\mathbf{v} = (v_1, v_2, v_3)$ and $\mathbf{w} = (w_1, w_2, w_3)$ in $\mathbb{R}^3$,
the \emph{cross product} is
\[
  \mathbf{v} \times \mathbf{w}
  = \begin{vmatrix} \mathbf{e}_1 & \mathbf{e}_2 & \mathbf{e}_3 \\
                    v_1 & v_2 & v_3 \\
                    w_1 & w_2 & w_3
    \end{vmatrix}
  = (v_2 w_3 - v_3 w_2,\;\; v_3 w_1 - v_1 w_3,\;\; v_1 w_2 - v_2 w_1).
\]
\end{tcolorbox}

\begin{tcolorbox}[colback=thmbox, colframe=thmborder, arc=2pt,
  left=6pt, right=6pt, top=4pt, bottom=4pt,
  title={\small\textbf{Theorem (Properties of the Cross Product)}},
  fonttitle=\small\bfseries]
For $\mathbf{v}, \mathbf{w}, \mathbf{x} \in \mathbb{R}^3$ and $c \in \mathbb{R}$:
\begin{enumerate}
  \item $\mathbf{v} \times \mathbf{w} = -(\mathbf{w} \times \mathbf{v})$ \hfill (anticommutativity)
  \item $\mathbf{v} \times (\mathbf{w} + \mathbf{x}) = \mathbf{v} \times \mathbf{w} + \mathbf{v} \times \mathbf{x}$ \hfill (distributivity)
  \item $\mathbf{v} \times \mathbf{v} = \mathbf{0}$
  \item $(c\mathbf{v}) \times \mathbf{w} = c(\mathbf{v} \times \mathbf{w})$
  \item $\mathbf{v} \cdot (\mathbf{v} \times \mathbf{w}) = 0$ and $\mathbf{w} \cdot (\mathbf{v} \times \mathbf{w}) = 0$ \hfill (orthogonality)
  \item $\|\mathbf{v} \times \mathbf{w}\| = \|\mathbf{v}\|\,\|\mathbf{w}\|\sin\theta$ \hfill (magnitude = parallelogram area)
\end{enumerate}
\end{tcolorbox}

\begin{remark*}[NURBS connection: surface normal vectors]
The unit normal to a parametric surface at a point $(u,v)$ is
\[
  N(u,v) = \frac{S_u \times S_v}{\|S_u \times S_v\|},
\]
where $S_u$ and $S_v$ are the partial derivative vectors of the surface
with respect to the parameters.
The cross product gives the direction orthogonal to both tangent directions;
dividing by the norm produces the unit normal.
This appears in Piegl \& Tiller §1.1 and is used throughout whenever
normal vectors are needed (shading, offset surfaces, etc.).
\end{remark*}

\begin{remark*}[Mnemonic for the cross product]
Write the entries of $\mathbf{v}$ and $\mathbf{w}$ each twice in a $2 \times 6$ array.
Ignore the first and last columns.
Draw three ``X'' shapes between columns 2–3, 3–4, and 4–5.
Each ``X'' contributes one entry of $\mathbf{v} \times \mathbf{w}$: add along forward
diagonals, subtract along backward diagonals.
\end{remark*}
